\documentclass[11pt,a4paper]{article}

\usepackage[T1]{fontenc}
\usepackage[utf8]{inputenc}
\usepackage[french]{babel}
\usepackage{lmodern}
\usepackage{amsmath,amssymb}
\usepackage{graphicx}
\usepackage{booktabs}
\usepackage{microtype}
\usepackage[margin=2.4cm]{geometry}
\usepackage[hidelinks]{hyperref}
\usepackage{natbib}
\bibliographystyle{plainnat}

\graphicspath{{figures/}}

\title{\textbf{Mesurer la diversité qu'un fil algorithmique expose réellement}\\[2mm]
       \large Une norme, ses attaques, et ce qu'il faut pour la vérifier}
\author{Sylvain Geffroy}
\date{Août 2026 --- \emph{document de travail ouvert à la relecture}}

\begin{document}
\maketitle

\begin{abstract}
\noindent
Savoir ce qu'un algorithme de recommandation expose réellement suppose une grandeur qu'on
puisse constater sans accéder au code de la plateforme. Cette note en propose une ---
l'\textbf{Indice de Diversité Exposée} (IDE), l'entropie des contenus servis projetés sur un
catalogue de points de vue déclaré, pondérée par l'attention qu'appelle chaque rang --- puis la
soumet aux attaques qu'elle doit supporter pour valoir comme mesure.

Trois de ces attaques réussissent contre les formulations naïves, et chacune impose une
correction. Un plancher portant sur les \emph{étiquettes} d'un fil est saturé à coût nul :
une plateforme capable de dissocier l'étiquette du contenu obtient la note maximale pour une
diversité de contenu strictement nulle. L'entropie quadratique de Rao, premier remplaçant
envisagé, a un optimum sous contrainte \emph{bimodal} : elle prescrirait la polarisation
qu'elle prétend mesurer. Et toute mesure aveugle au rang se laisse satisfaire en
\emph{enterrant} les contenus divergents --- une plateforme certifiée à $0{,}70$ n'expose
que $0{,}36$ de diversité réelle, et fermer l'échappatoire double le coût d'engagement.

Vérifier une telle norme sur données réelles exige de connaître l'\emph{exposition} : nous
donnons trois contrôles à passer sur un journal de recommandation --- l'ordre y dit-il quelque
chose, y a-t-il de quoi estimer, et sous quelle forme --- et les appliquons à trois journaux
publics. MIND ne conserve pas l'ordre servi, et l'estimateur usuel y produit cinq sévérités
incompatibles ; Baidu-ULTR le conserve ; l'Open Bandit Dataset permet la seule confrontation de
ce travail à une vérité terrain, où la valeur d'une politique jamais déployée est retrouvée à
$2{,}5\,\%$ contre $31{,}6\,\%$ d'erreur pour l'estimation naïve.

La conclusion principale de cette partie est cependant négative pour la méthode que nous venions
de construire. Baidu-ULTR publie une colonne de temps d'affichage qui permet de \textbf{mesurer}
l'exposition au lieu de l'estimer à travers les clics : la sévérité y vaut alors
$0{,}88 \pm 0{,}05$ sur $143$ documents, et non $1{,}09 \pm 0{,}09$ sur $55$ --- l'estimation
surestime la décroissance de $23\,\%$, parce qu'un clic mélange l'examen et l'attrait. La mesure
directe réfute en outre le modèle à cascade sur ce journal, par deux voies indépendantes, et
montre que la loi $R^{-\eta}$ est le pire des trois ajustements sur la courbe observée. Nous avions enfin conclu
qu'un format d'affichage enrichi placé au-dessus retire $8{,}4$ points d'examen à ce qui suit, et
que la pondération de rang était donc une propriété de la \textbf{page servie} plutôt que du rang.
Une mesure appariée retire cette conclusion : à \emph{contenu} tenu fixe, il n'en reste que $6\,\%$
(rapport de risques $0{,}940$, IC$_{95}$ $[0{,}916 ; 0{,}964]$), les deux tiers de l'effet publié
étant une composition de requêtes. L'incidence sur l'indice est de $0{,}0025$, contre $0{,}0350$
pour la convention $1/R$ : la loi de rang est donc \textbf{transportable}, et ce qui menace la
mesure n'est pas la page mais de n'avoir pas mesuré.

Un quatrième journal, EB-NeRD, porte enfin les deux colonnes attendues --- la liste servie et
une rubrique déclarée. Le premier contrôle dit que son \textbf{ordre enregistré ne contient pas
le rang} ($z = +1{,}05$, alors que le journal aurait détecté une sévérité cent trente-quatre fois
plus faible). L'indice aveugle au rang y est donc mesuré pour la première fois sur des fils réels
--- $0{,}50$ par journée-utilisateur, soit $5{,}1$ rubriques également servies sur $26$ --- et
l'indice exposé seulement \emph{encadré}, à $0{,}104$ près : à un plancher de $0{,}40$, un tiers
des fils y contreviennent quelle que soit leur mise en ordre, et six sur dix restent indécidables.
Ce qu'il faut réclamer d'un journal n'est donc pas le rang, mais le rang \textbf{vérifiable} ---
et le test d'échangeabilité est ce qui le vérifie. Le catalogue, enfin, décide du niveau bien plus
que le fil lui-même : le même corpus vaut $0{,}50$ sur vingt-six rubriques et $0{,}92$ sur trois,
tandis que l'ordre entre lecteurs tient jusqu'à six modalités puis s'effondre.

L'indice est issu d'un modèle thermodynamique dont l'audit a réfuté chaque affirmation
propre, et qui ne figure plus ici : l'analogie avec la décohérence
quantique qui a lancé ce travail ne survit à aucun de ses transferts. L'ensemble du code,
des tests et des figures est reproductible en conteneur.
\end{abstract}

\part{La mesure}
\label{part:mesure}

\section{Le problème métrologique}
\label{sec:probleme}

Réguler un algorithme de recommandation suppose une grandeur constatable. Compter les
contenus problématiques ne convient pas : l'indicateur est retardé, manipulable, et il
demande de trancher sur le fond de chaque contenu. La grandeur retenue ici est d'une autre
nature --- elle porte sur la \emph{forme} de l'exposition, non sur le contenu des messages :

\begin{quote}
Quelle part de l'attention servie à un lecteur va à chacun des points de vue que la
plateforme aurait pu lui servir ?
\end{quote}

Cette question a trois qualités : elle se mesure depuis l'extérieur, elle ne demande aucun
jugement sur la véracité d'un contenu, et elle admet un seuil exprimé en pourcentage, donc
transposable d'une plateforme à l'autre.

Elle a aussi trois pièges, et cette note est organisée autour d'eux.

\begin{enumerate}
\item \textbf{Une étiquette n'est pas un contenu.} Mesurer la diversité des \emph{catégories
      annoncées} laisse à la plateforme le choix de ce qui est annoncé.
\item \textbf{Une composition n'est pas une exposition.} Un fil peut contenir la diversité
      exigée et la placer là où personne ne regarde.
\item \textbf{Un clic enregistré n'est pas une préférence.} Les clics observés ont été
      produits par le classement de la plateforme, et non par celui qu'on évalue.
\end{enumerate}

Chacun de ces pièges est mesurable, et chacun a été mesuré. Ce qui subsiste après les trois
est la définition retenue.

\section{L'indice de diversité exposée}
\label{sec:indice}

Soit un fil de $n$ positions, un catalogue de référence de $k$ points de vue déclaré par le
déclarant, et une remise d'attention $w_R$ associée au rang $R$. La distribution
\emph{exposée} est
\begin{equation}
q_i = \frac{\sum_{R=1}^{n} w_R \, \mathbb{1}\!\left[\text{le contenu servi au rang } R
      \text{ tombe dans le bac } i\right]}{\sum_{R=1}^{n} w_R},
\label{eq:distribution}
\end{equation}
et l'indice est son entropie de Shannon normalisée :
\begin{equation}
\mathrm{IDE} = \frac{H(q)}{\log_2 k}, \qquad H(q) = -\sum_{i=1}^{k} q_i \log_2 q_i .
\label{eq:ide}
\end{equation}

Trois choix distinguent \eqref{eq:distribution} de l'entropie qu'on écrirait spontanément,
et chacun est la conséquence d'une attaque qui a réussi (\S\ref{sec:adverse}).

\begin{center}
\begin{tabular}{@{}p{0.30\linewidth}p{0.30\linewidth}p{0.32\linewidth}@{}}
\toprule
\textbf{Choix} & \textbf{Raison} & \textbf{Sans lui} \\
\midrule
$q$ porte sur les contenus servis, projetés sur les bacs du catalogue & une étiquette se
choisit, un contenu se constate & l'indice vaut $1{,}000$ pour une diversité de contenu
nulle, à coût nul \\[2mm]
chaque rang est pondéré par $w_R$ & un lecteur consulte le premier élément bien plus souvent
que le dernier & la norme se satisfait en enterrant : $0{,}70$ certifié, $0{,}36$ exposé \\[2mm]
le dénominateur $\log_2 k$ vient du catalogue déclaré & comparer deux plateformes exige la
même unité & un fil fermé ne présente qu'une modalité, le dénominateur dégénère et l'indice
flatte \\
\bottomrule
\end{tabular}
\end{center}

\paragraph{Remise d'attention.} Les résultats de cette note emploient $w_R = 1/R$, la remise
du rang réciproque, usuelle en évaluation de classement. Le choix de $w$ appartient au
déclarant au même titre que le catalogue ; ce qui compte ici est qu'il soit \emph{déclaré},
car c'est lui qui rend la norme non contournable par l'ordre.

\paragraph{Grandeur de contrôle associée.} L'écart entre la valeur de \eqref{eq:ide} calculée
sur la composition du fil ($w_R$ constant) et sur sa distribution exposée ($w_R = 1/R$)
vaut zéro pour une plateforme qui ne relègue pas, et croît avec l'enterrement. C'est la
seule grandeur de ce travail qui compare une mesure \emph{à elle-même}, et donc la seule qui
se seuille sans référence extérieure.

\paragraph{Comment publier le chiffre.} Une entropie normalisée n'est pas une diversité : elle
n'est pas linéaire en ce qu'on entend par « deux fois plus divers ». Sa conversion en
\textbf{nombre effectif de points de vue}, ${}^{1}\!D = k^{\mathrm{IDE}}$, l'est
\citep{jost2006} : un plancher de $0{,}70$ sur un catalogue de quatre vaut $2{,}6$ points de vue
effectifs, et les $0{,}44$ réellement exposés par le fil enterrant en valent $1{,}9$.

\paragraph{Protocole.} L'unité est le contenu \emph{servi}, non le contenu consulté : c'est
la plateforme qui est auditée, pas le lecteur. L'indice se calcule par lecteur sur une
fenêtre glissante, puis s'agrège ; la grandeur réglementaire est la \emph{part de la
population sous le seuil}, non la moyenne, car une moyenne satisfaisante peut masquer une
minorité entièrement enfermée.

\section{L'épreuve adverse}
\label{sec:adverse}

Une norme se juge à ce qu'elle contraint un adversaire rationnel à faire. Nous plaçons donc
la plateforme dans le rôle qu'elle occupera : maximiser l'engagement \emph{sous} la
contrainte, et non s'y conformer de bonne foi.

\paragraph{Méthode.} Le fil compte $n$ positions à remplir depuis un catalogue de $k$ points
de vue, soit $k^n$ fils possibles. Pour $n = 8$ et $k = 4$, les $65\,536$ fils sont
\emph{tous énumérés} : l'optimum obtenu est exact, non le résultat d'une heuristique. Ce
n'est pas un luxe. Il s'agit de résultats négatifs --- des normes qui échouent --- et un
optimum manqué par un solveur produirait exactement l'apparence d'une norme qui tient.

\subsection{Une norme sur les étiquettes est saturable à coût nul}

Si la plateforme contrôle l'étiquette indépendamment du contenu, l'entropie des étiquettes
cesse de contraindre l'entropie des contenus. L'optimisation le confirme sans nuance : sous
découplage complet, la plateforme obtient un indice de $1{,}000$ --- la note maximale --- pour
une diversité de contenu strictement nulle, sans céder un point d'engagement. Il n'est même
pas nécessaire d'aller jusque-là : à mi-découplage, la contrainte n'a plus que $36\,\%$ de sa
force.

La conséquence est directe : \textbf{la mesure doit porter sur les contenus servis}. Elle
appelle un diagnostic complémentaire, l'\emph{excès de signature} --- l'écart entre deux
indices calculés sur le même fil, rapporté à ce qu'un catalogue honnête afficherait au même
indice --- qui vaut zéro pour une plateforme honnête et croît avec le découplage.

\subsection{Le premier remplaçant prescrivait la polarisation}

L'entropie quadratique de Rao \citep{rao1982} --- la distance moyenne entre deux contenus
tirés du fil, connue en recommandation sous le nom d'\emph{intra-list distance} --- semblait
répondre à l'objection : elle voit les contenus, non les étiquettes. Son optimum sous
contrainte d'engagement est pourtant \textbf{bimodal} : à diversité imposée, la façon la
moins coûteuse de la fournir est de servir deux blocs extrêmes et de vider le centre. Une
norme fondée sur elle prescrirait la polarisation qu'elle prétend mesurer. Ce défaut est
connu de la littérature sur la diversification \citep{ohsaka2023} ; il est ici retrouvé
comme conséquence directe de l'optimisation sous contrainte.

Le plancher retenu porte donc sur une \emph{entropie}, calculée sur les contenus servis
projetés sur les bacs du catalogue --- mesure dont l'optimum sous contrainte est étalé et
non bimodal.

\subsection{Une norme aveugle au rang se contourne par l'enterrement}

Les mesures ci-dessus portent sur la \emph{composition} du fil. Reprises sur des fils
\emph{ordonnés}, les quatre mesures de diversité candidates se laissent toutes contourner de
la même façon : servir les contenus divergents aux derniers rangs, où l'attention ne va plus.

\begin{center}
\begin{tabular}{@{}lrrrr@{}}
\toprule
\textbf{Mesure} & \textbf{Coût} & \textbf{Affiché} & \textbf{Exposé} & \textbf{Coût si conscient du rang} \\
\midrule
Entropie de Rao & $8{,}2\,\%$ & $0{,}750$ & $\mathbf{0{,}355}$ & $18{,}9\,\%$ \\
Entropie de position & $10{,}7\,\%$ & $0{,}774$ & $\mathbf{0{,}443}$ & $20{,}9\,\%$ \\
Gaussian ILD & \multicolumn{4}{l}{plancher inatteignable} \\
Proximité à la cible & $5{,}9\,\%$ & $0{,}750$ & $\mathbf{0{,}628}$ & $10{,}6\,\%$ \\
\bottomrule
\end{tabular}
\end{center}

Le fil optimal a toujours la même forme : les contenus du point de vue préféré en tête, les
divergents relégués. \textbf{Une plateforme certifiée à $0{,}70$ n'expose que $0{,}36$.} Un
plancher conscient du rang ferme l'échappatoire et \emph{double} le coût d'engagement : une
norme qui ne coûterait pas davantage ne fermerait rien.

L'écart croît avec l'exigence --- $0{,}262$ à plancher $0{,}40$, $0{,}395$ à plancher
$0{,}70$ pour l'entropie de Rao : plus la norme aveugle demande de diversité, plus il devient
rentable de l'enterrer.

\begin{figure}[t]
\centering
\includegraphics[width=\linewidth]{fig15_rang_adverse.png}
\caption{Diversité affichée contre diversité exposée sous plancher aveugle au rang ; écart
selon l'exigence du plancher ; sévérité du biais de position retrouvée par effets fixes, et
seuil d'exploration au-delà duquel elle cesse de l'être ; coût d'un $\eta$ posé au jugé.}
\label{fig:adverse}
\end{figure}

\section{Mesurer l'exposition}
\label{sec:exposition}

Tout ce qui précède suppose de savoir quelle attention va à quel rang. Le modèle usuel pose
une probabilité d'examen décroissante avec le rang \citep{joachims2017},
\begin{equation}
e(R) = R^{-\eta},
\label{eq:examen}
\end{equation}
dont la sévérité $\eta$ était, dans la première version de ce travail, \emph{posée}. Elle
s'estime.

\subsection{Estimation par effets fixes de contenu}

Sous \eqref{eq:examen}, la probabilité de clic se factorise en $P(\text{clic} \mid i, R) =
g(i)\, R^{-\eta}$, donc
\begin{equation}
\log \mathrm{CTR}(i, R) = \log g(i) - \eta \log R .
\label{eq:effetsfixes}
\end{equation}
La pertinence $g(i)$ y est un effet fixe de contenu : on ne l'estime pas, on l'élimine par
centrage intra-contenu. Ce qui subsiste est la seule variation qui identifie $\eta$ --- celle
d'un \emph{même contenu vu à des rangs différents}. C'est la forme la plus simple de la
récolte d'interventions \citep{agarwal2019} : elle n'exige aucune expérience, seulement que
la plateforme n'ait pas toujours classé les mêmes contenus aux mêmes places.

Sur données simulées, l'estimateur retrouve $\eta$ de $0{,}4$ à $1{,}6$. À politique
déterministe, il \emph{refuse} de répondre plutôt que d'inventer un chiffre : aucun contenu
ne change de rang, la sévérité n'est pas dans les données.

Poser $\eta$ au lieu de l'estimer n'est pas anodin : sur un cas où le coût réel d'un filtre
de diversité vaut $6{,}6\,\%$, un $\eta$ supposé de $0{,}5$ le sous-estime de $59\,\%$ et un
$\eta$ supposé de $2{,}0$ le surestime de $179\,\%$ --- l'ordre de grandeur même du biais
qu'on prétendait corriger.

\subsection{Le contrôle qui doit précéder : l'échangeabilité}

Le refus de répondre décrit ci-dessus attrape le cas visible --- aucune variation de rang ---
et laisse passer le cas invisible : une variation \emph{abondante et artificielle}. Il faut
donc un contrôle qui ne compte pas la variation mais l'\emph{informativité} de l'ordre.

Pour un fil de longueur $L$ portant $k$ clics, soit $u_R = (R - 1/2)/L$ le rang normalisé.
Sous l'hypothèse que les clics sont indifférents à la position \emph{à fil donné}, les
positions cliquées sont un tirage sans remise parmi les $L$ positions, d'espérance et de
variance connues exactement :
\begin{equation}
\mathbb{E}\Big[\sum_{\text{clics}} u_R\Big] = k\,\bar{u}, \qquad
\mathbb{V}\Big[\sum_{\text{clics}} u_R\Big] = \frac{k(L-k)}{L-1}\,\sigma^2_u .
\label{eq:echangeabilite}
\end{equation}
Le conditionnement au fil rend le test immunisé au mélange des longueurs, à la qualité
moyenne des contenus d'un fil et à l'appétit de clic de son lecteur. Un biais de position
concentre les clics en haut : il rend l'écart réduit \emph{négatif}, et le signe du rejet est
donc lui-même une vérification.

Un test qui ne rejette rien ne dit rien tant qu'on ignore ce qu'il rejetterait. L'étalonnage
se fait par simulation à structure de fils identique : sur celle de MIND, le test détecte
$\eta = 0{,}02$ à douze écarts-types, et la sévérité minimale détectable vaut
$\eta \approx 0{,}004$.

\section{Évaluer un réordonnancement}
\label{sec:evaluation}

Mesurer un fil est une chose ; estimer ce que coûterait un autre classement en est une autre.
Les clics disponibles ont été produits par la politique de la plateforme, non par celle qu'on
évalue.

\paragraph{Le \emph{replay} est faux, et d'un montant considérable.} Rejouer les clics
enregistrés sur un fil réordonné se trompe de $201\,\%$ en médiane, jusqu'à $851\,\%$, et le
sens de l'erreur n'est pas garanti : sur soixante tirages à configuration identique, il
surestime dans cinquante-six cas et sous-estime dans quatre. Un chiffre naïf n'est donc même
pas une borne supérieure.

\paragraph{La correction et son prix.} La pondération par l'inverse de la propension
(\emph{IPS}) est sans biais dès lors que la politique d'enregistrement est connue et donne
une chance à toute action que la politique cible pourrait choisir. Son prix est la variance,
d'où les variantes usuelles --- auto-normalisation \citep{swaminathan2015}, plafonnement des
poids --- et le diagnostic qui doit accompagner tout chiffre : la \textbf{taille
d'échantillon effective}, qui dit combien d'observations portent réellement l'estimation.

\paragraph{Une confrontation à la vérité terrain.} L'Open Bandit Dataset
\citep{saito2020} contient deux seaux servis en parallèle par deux politiques différentes,
et publie les propensions vraies. On peut donc estimer la valeur de la politique uniforme
\emph{sur les seules données d'une autre politique}, puis la comparer à sa valeur mesurée là
où elle a réellement servi.

\begin{center}
\begin{tabular}{@{}lrr@{}}
\toprule
\textbf{Estimateur} & \textbf{Valeur} & \textbf{Écart à la vérité} \\
\midrule
Vérité terrain (mesurée, $452\,949$ impressions) & $0{,}005124$ & --- \\
Naïf (taux de clic observé) & $0{,}006743$ & $+31{,}6\,\%$ \\
IPS & $0{,}005253$ & $\mathbf{+2{,}5\,\%}$ \\
Auto-normalisé & $0{,}004768$ & $-7{,}0\,\%$ \\
IPS plafonné à $10$ & $0{,}004473$ & $-12{,}7\,\%$ \\
\bottomrule
\end{tabular}
\end{center}

La correction fonctionne. Le diagnostic interdit toutefois d'en tirer gloire : la taille
d'échantillon effective vaut $1\,513$ pour $4\,077\,727$ impressions, soit $0{,}04\,\%$.
L'estimation est sans biais et repose sur l'équivalent de mille cinq cents observations ;
que l'écart tombe à $2{,}5\,\%$ tient autant à la chance qu'à la méthode. C'est le cas réel
qui justifie la règle : publier la taille effective à côté du chiffre.

\section{Trois journaux publics}
\label{sec:journaux}

Les instruments des \S\ref{sec:exposition} et \S\ref{sec:evaluation} décident si un jeu de
données permet la mesure. Appliqués à trois journaux publics, ils donnent trois réponses
différentes.

\subsection{MIND : un ordre qui n'en est pas un}

Le jeu de référence de la recommandation d'actualité \citep{wu2020} présente une couverture
de rang idéale : un contenu médian y apparaît à seize positions distinctes. Le test
\eqref{eq:echangeabilite} n'y détecte pourtant rien --- $z = +0{,}12$ ($p = 0{,}91$) sur
$156\,965$ fils, répliqué à $z = +0{,}28$ sur le second découpage --- dans un jeu où il
détecterait $\eta = 0{,}02$ à douze écarts-types. L'ordre enregistré est indiscernable d'un
mélange.

La courbe de taux de clic par position y décroît pourtant de $0{,}108$ à $0{,}038$ et
s'ajuste à $\hat\eta = 0{,}39$ : c'est un \emph{artefact de composition}, les positions
profondes n'existant que dans les fils longs, dont le taux de clic par contenu est
mécaniquement plus faible. À longueur de fil fixée, la pente est nulle et change de signe
d'une longueur à l'autre.

Enfin l'estimateur \eqref{eq:effetsfixes} accepte ce jeu, se déclare identifiable, et rend
une sévérité qui est une fonction croissante d'un simple seuil de nuisance : $-0{,}131$ à
cinq impressions, $+0{,}053$ à vingt, $+0{,}255$ à cent, toutes assorties d'une erreur type
inférieure à $0{,}007$. \textbf{Cinq chiffres confiants et incompatibles tirés du même jeu.}

Le mélange ne débiaise pas les clics : ceux-ci ont été produits sous l'ordre réel et en
portent le biais en entier. Ce que le mélange retire est le \emph{rang}, seul régresseur qui
permettrait d'en tenir compte. Sur un journal simulé de sévérité vraie $1{,}00$, l'estimation
passe de $1{,}003$ à $-0{,}003$ après mélange, à clics identiques ; l'évaluation d'un
réordonnancement y devient vide par construction.

\subsection{Baidu-ULTR : le contrôle positif}

Sur une tranche de $524\,164$ documents servis en $64\,200$ sessions de recherche
\citep{zou2022}, le même test rejette à $z = -205{,}7$ ($p < 10^{-12}$) et du côté que la
théorie prescrit. La sévérité y vaut $\hat\eta = 1{,}10 \pm 0{,}09$ par effets fixes de
document, contre $1{,}49$ par ajustement agrégé --- l'écart étant le confondant de qualité,
la plateforme plaçant les meilleurs documents en tête. La couverture reste mince : $55$
documents portent l'estimation, sur $444\,709$, parce qu'une même URL réapparaît rarement.

\subsection{Open Bandit Dataset : la position mesurée sans modèle}

Le seau aléatoire affecte les contenus aux positions au hasard : l'écart de taux de clic
entre positions y est un effet \emph{causal}, sans modèle d'examen à poser. Sur un bandeau de
trois vignettes horizontales, il vaut $\hat\eta = 0{,}037$ (campagne « all », $1{,}37$
million d'impressions) à $\hat\eta = 0{,}105$ (campagne « men »).

\textbf{$\eta$ est donc une propriété de la surface, non une constante} : un ordre de grandeur
sépare une page de résultats verticale d'un bandeau de trois vignettes. Transporter la valeur
de l'une à l'autre est exactement l'erreur dont \S\ref{sec:exposition} chiffre le coût.

\begin{figure}[t]
\centering
\includegraphics[width=\linewidth]{fig17_rang_servi.png}
\caption{Ce qu'un rang réellement servi fait à la courbe de taux de clic ; le test
d'échangeabilité sur MIND et sur Baidu-ULTR ; l'effet de position mesuré sans modèle sur le
seau aléatoire de l'Open Bandit Dataset ; l'estimateur contrefactuel jugé contre la valeur
qu'il estime.}
\label{fig:journaux}
\end{figure}

\subsection{Ce que ces trois journaux ne permettent pas}

Aucun ne porte à la fois le \emph{rang servi} et une \emph{étiquette de point de vue
interprétable}. MIND a les catégories éditoriales sans le rang ; Baidu-ULTR a le rang sans
étiquette ; l'Open Bandit Dataset a le rang, la propension et trois attributs catégoriels,
mais anonymisés --- et une diversité calculée sur des hachages ne dit rien, le catalogue de
référence étant une déclaration politique.

L'indice de la \S\ref{sec:indice} n'a donc jamais été mesuré sur un fil réel. Ce n'est pas
une lacune de méthode : c'est une lacune de donnée.

\section{Limites}
\label{sec:limites}

Ces limites sont énoncées ici plutôt que laissées aux relecteurs.

\subsection{Sur l'indice et la norme}

\paragraph{La forme retenue n'a jamais été mesurée sur un fil réel.} Elle est définie
\eqref{eq:ide}, son coût est chiffré par énumération exhaustive, et c'est tout. La lacune est
de \emph{donnée}, non de méthode (\S\ref{sec:journaux}) : il y faudrait un journal qui
faudrait pour la combler.

\paragraph{Le niveau du plancher n'est pas déductible de la mesure.} Ce travail établit la
forme de la norme et son prix, pas sa valeur. Fixer $0{,}60$ plutôt que $0{,}40$ est une
décision politique qu'aucun calcul présenté ici ne tranche.

\paragraph{La discrétisation en points de vue est un choix politique.} Qui définit les
modalités définit l'indice. La mesure par divergence à un catalogue déclaré rend ce choix
explicite au lieu de l'enfouir ; elle ne le supprime pas.

\paragraph{L'indice reste manipulable au-delà de ce qu'une mesure automatique détecte.} Les
attaques de \S\ref{sec:adverse} sont fermées, mais une plateforme peut encore servir des
contenus formellement divergents et substantiellement vides --- de la diversité de bac sans
diversité d'argument. Toute métrique imposée devient une cible.

\paragraph{Un plancher est une contrainte sur ce que les citoyens voient.} C'est défendable,
mais c'est une intervention sur le débat public, et la présenter comme une simple mesure
technique serait malhonnête. La mesure sur des fils individuels soulève en outre une question
de vie privée à laquelle cette note ne répond pas.

\subsection{Sur les mesures d'exposition}

\paragraph{La forme $e(R) = R^{-\eta}$ est posée.} Seule sa sévérité est estimée. Un modèle
d'examen à cascade, ou dépendant du contenu, produirait une autre dépendance au rang --- que
le test \eqref{eq:echangeabilite} détecterait tout aussi bien, puisqu'il ne teste que
l'indépendance, mais que la sévérité minimale détectable ne résumerait plus.

\paragraph{L'énumération exhaustive borne la taille des fils.} Huit positions sur quatre
points de vue. Rien n'assure que l'enterrement se comporte ainsi sur un fil de cinquante
items, et l'affirmer demanderait une optimisation approchée dont il faudrait établir qu'elle
n'invente pas le résultat.

\paragraph{Les jeux mesurés sont des tranches.} Une tranche de Baidu-ULTR n'est pas
Baidu-ULTR ; l'Open Bandit Dataset mesure une recommandation de mode sur trois vignettes
horizontales, et le transport de ses chiffres vers un fil d'actualité est précisément
l'opération que \S\ref{sec:journaux} déconseille.

\paragraph{La confrontation à la vérité terrain porte sur une politique cible uniforme}, la
seule dont ce jeu contienne un seau. Rien n'assure que la pondération d'importance se comporte
aussi bien pour une cible plus éloignée de la politique d'enregistrement --- au contraire, la
taille d'échantillon effective y serait plus faible encore.



\subsection{Les deux hypothèses les plus profondes, éprouvées}
\label{sec:anglesmorts}

Deux hypothèses traversent toute cette note sans avoir été mises à l'épreuve : que l'examen suit
une loi de puissance, et que la pertinence est connue. Aucune ne l'est.

\paragraph{Le test d'échangeabilité survit à un changement de modèle de clic.} Sous un modèle
\emph{à cascade} \citep{craswell2008} --- le lecteur descend le fil, s'arrête dès qu'il a trouvé, et poursuit sinon
avec une probabilité $\gamma$ --- le test \eqref{eq:echangeabilite} rejette entre $z = -208$ et
$z = -241$ selon $\gamma$, soit \emph{plus fortement} que sur Baidu-ULTR. C'est le résultat qu'il
fallait obtenir : le test ne suppose rien de la forme de l'attention, il ne teste que
l'indépendance, et son verdict négatif sur MIND garde donc son sens sous n'importe quel modèle.

\paragraph{La loi de puissance, elle, ne survit pas.} Sous cascade, la décroissance de l'examen
est \emph{géométrique} et non polynomiale. La simulation rend ici la vérité terrain que les
données réelles ne donnent jamais --- les positions effectivement examinées --- et la
comparaison est sans appel : au douzième rang, la loi $R^{-\hat\eta}$ ajustée surestime
l'exposition réelle d'un facteur $40$ à $\gamma = 0{,}95$, $129$ à $\gamma = 0{,}85$ et $4\,110$
à $\gamma = 0{,}60$. L'ajustement paraît pourtant excellent, avec une erreur type de $0{,}04$.

La conséquence porte directement sur \S\ref{sec:journaux} : la sévérité
$\hat\eta = 1{,}10 \pm 0{,}09$ mesurée sur Baidu-ULTR n'a de sens \emph{que si} l'examen y suit
une loi de puissance --- et Baidu-ULTR est une page de résultats de recherche, terrain d'élection
du modèle à cascade. La sévérité doit donc être publiée \textbf{avec la forme supposée}, et cette
forme reste à tester. Le test d'échangeabilité dit si l'ordre porte de l'information ; il ne dit
pas sous quelle forme.

\paragraph{Et une pertinence estimée efface l'avantage des méthodes réglées.} La comparaison de
La comparaison de réordonnanceurs, retirée de cette note, supposait la pertinence connue. En faisant classer les méthodes sur une
pertinence \emph{perçue} et en les jugeant sur la vraie, leur ordre ne change pas --- le filtre et
MMR restent devant --- mais leur avantage sur le tirage au sort passe de $16{,}4$ points à bruit
nul à $5{,}8$ points pour une corrélation de rang de $0{,}50$, et \emph{s'inverse} en deçà de
$0{,}36$. La raison est mécanique : le tirage au sort est le seul réordonnanceur qui n'utilise pas
la pertinence, donc le seul dont le manque à gagner soit invariant au bruit.

À pertinence réaliste, l'écart entre réordonnanceurs compte donc bien moins que la qualité du
moteur de pertinence sous-jacent --- ce qui renforce, par une troisième voie, la conclusion que la
contribution de ce travail n'est pas son algorithme.


\subsection{Un troisième contrôle : la forme de l'examen}
\label{sec:testdeforme}

Le test d'échangeabilité dit si l'ordre d'un journal porte de l'information ; il ne dit pas sous
quelle forme, et \S\ref{sec:anglesmorts} a montré que la forme décide de tout. Un troisième
contrôle s'impose donc, et il repose sur une \emph{indépendance conditionnelle} plutôt que sur une
allure : sous un modèle de position, l'examen du rang $R$ ne dépend que de $R$, donc le clic y est
indépendant de ce qui s'est passé au-dessus ; sous cascade, un clic au-dessus supprime l'examen en
dessous.

Pour chaque cellule $s = (\text{contenu}, \text{rang})$, les $n_s$ impressions se répartissent en
$n_{1s}$ précédées d'un clic et $n_{0s}$ qui ne le sont pas, pour $k_s$ clics au total. Sous
l'hypothèse d'indépendance, ces clics se répartissent comme un tirage sans remise :
\begin{equation}
\mathbb{E}[a_s] = \frac{k_s n_{1s}}{n_s}, \qquad
\mathbb{V}[a_s] = \frac{k_s (n_s - k_s)\, n_{1s} n_{0s}}{n_s^2 (n_s - 1)} .
\label{eq:forme}
\end{equation}
C'est la statistique de Mantel-Haenszel, stratifiée par cellule ; le conditionnement élimine la
qualité du contenu.

\paragraph{Le contrôle fonctionne.} Il ne rejette jamais sous modèle de position ($|z| < 1$ pour
$\eta \in \{0 ; 0{,}5 ; 1 ; 2\}$) et rejette massivement sous cascade ($z = -97$ à $-307$).

\paragraph{Mais il ne tranche pas ce qu'il devait trancher.} Un lecteur qui cesse de cliquer une
fois servi --- tout en continuant de parcourir le fil --- produit la même signature : $z = -144$
sous modèle de position \emph{pur} avec un budget de un clic, contre $-179$ sous cascade. Rien
dans les clics ne les distingue, et c'est attendu : dans les deux cas le journal montre qu'après
un clic il n'y a plus rien. La différence est pourtant décisive, puisque sous budget le contenu
placé sous un clic est examiné, et sous cascade il ne l'est pas.

\paragraph{Sur Baidu-ULTR, aucune signature de cascade.} L'écart réduit y est \emph{positif} ---
$+5{,}6$ au seuil le plus permissif, $+0{,}5$ à $+0{,}6$ aux plus stricts --- alors qu'une cascade
le rendrait négatif. Le signe est celui du confondant d'hétérogénéité, que le test annonçait comme
sa limite. Trois lectures restent ouvertes : l'examen y est proche d'un modèle de position ; une
cascade existe mais le confondant la masque ; ou la couverture est trop mince, $108$ cellules au
seuil le plus strict.

\paragraph{Une erreur de protocole, rattrapée.} Restreindre le journal aux sessions à plusieurs
clics --- où un budget de un est exclu par construction --- semblait la façon évidente de séparer
les deux modèles, et donne sur Baidu-ULTR $z = -8{,}2$ ($p = 2 \times 10^{-16}$). Le nombre de
clics d'un fil est pourtant un \emph{collider} de ses clics individuels : conditionner dessus
induit entre eux une dépendance négative, et fabrique donc la signature recherchée. Sur un journal
simulé sous modèle de position pur, la même restriction fait passer $z$ de $+0{,}74$ à $-45{,}7$.
Ce qui était mesuré n'était pas une propriété du journal mais la sélection opérée.

\paragraph{Ce qui trancherait.} Une mesure de l'\emph{examen} plutôt que du clic --- temps
d'affichage, profondeur de défilement, contenu sorti de l'écran. Baidu-ULTR publie ces colonnes ;
ce travail ne les a pas lues.


\subsection{L'exposition mesurée plutôt qu'estimée}
\label{sec:expositionmesuree}

Baidu-ULTR publie une colonne \texttt{displayed\_time} --- le temps d'affichage de chaque document
--- que ce travail n'avait pas lue. Elle mesure directement ce qui a été \emph{montré}, là où tout
ce qui précède n'observait que ce qui avait été \emph{cliqué}.

\paragraph{Elle lève l'ambiguïté de \S\ref{sec:testdeforme}.} Sur les clics, cascade et budget
sont confondus ; sur l'examen, la séparation est totale : $z = -158$ à $-412$ sous cascade contre
$z = -1{,}03$ sous budget, quel que soit le budget. Appliquée à Baidu-ULTR, la mesure \textbf{réfute
la cascade} : l'examen sous un clic y est \emph{plus} fréquent, pas moins --- $0{,}221$ contre
$0{,}136$ au neuvième rang, $z = +8{,}4$ à $+10{,}9$ selon le seuil. C'est le confondant
d'hétérogénéité, mesuré directement.

\paragraph{Elle révise le chiffre publié.} La sévérité de l'exposition, mesurée par effets fixes de
document sur l'affichage, vaut $0{,}882 \pm 0{,}046$ sur \textbf{143} documents, contre
$1{,}085 \pm 0{,}093$ sur $55$ pour l'estimation par les clics. L'estimation \textbf{surestime la
décroissance de $23\,\%$}, et l'écart n'est pas une imprécision mais un confondant identifié : un
clic est le produit de l'examen et de l'attrait, et l'attrait décroît lui aussi avec le rang.

\paragraph{Elle met en cause la forme employée partout.} Ajustées sur la courbe \emph{mesurée},
les trois formes candidates donnent $R^2 = 0{,}72$ pour la loi de puissance --- avec $42\,\%$
d'écart maximal ---, $0{,}96$ pour une décroissance géométrique et $0{,}99$ pour un modèle à
\emph{pli d'écran}. La loi en $R^{-\eta}$ est la pire des trois, et elle se trompe surtout là où la
norme se joue : au deuxième rang, elle prédit $0{,}47$ quand la mesure donne $0{,}88$.

\paragraph{Réserve.} Un temps d'affichage positif atteste que le document a été montré, non qu'il
ait été regardé. La mesure surestime donc l'exposition, dans une proportion que ce travail ne peut
pas chiffrer.


\subsection{Le format, le retour, et une explication retirée}
\label{sec:formatretour}

Deux colonnes de Baidu-ULTR restaient inexploitées : \texttt{slipoff\_count\_after\_click}, qui
compte les documents sortis de l'écran après un clic, et \texttt{media\_type}, qui distingue les
formats d'affichage.

\paragraph{La cascade réfutée une seconde fois.} La première colonne est non nulle pour
$43{,}7\,\%$ des lignes cliquées et $0{,}5\,\%$ des autres : après un clic, le journal enregistre
explicitement que des documents ont défilé, donc que le lecteur est revenu. Sous cascade stricte,
elle serait toujours nulle. La part croît d'ailleurs avec la profondeur du clic --- $0{,}37$ au
premier rang, $0{,}59$ au sixième. La réfutation de \S\ref{sec:expositionmesuree}, obtenue par un
test statistique, est ainsi confirmée par un fait consigné, ce qui est une voie indépendante.

\paragraph{Le format modifie l'exposition, et pas par la géométrie.} Un format enrichi ---
encadré de réponse, image, tableau --- est plus haut : $273$ pixels en médiane contre $192$. On
attendrait donc un effet de repoussement. À rang \emph{et} à hauteur cumulée comparables, un
format enrichi au-dessus retire pourtant $8{,}4$ points d'examen à ce qui suit, médiane sur $21$
strates toutes de même signe. L'effet ne passe donc pas par la place occupée, et l'explication qui
reste est celle de la \emph{satisfaction} : un encadré répond à la question, et le lecteur cesse de
descendre.

La conséquence porte sur toute mesure fondée sur le rang seul, y compris la nôtre : l'exposition
au rang $R$ dépend du \textbf{format de ce qui est au-dessus}, et aucune loi $e(R)$ ne peut la
représenter puisqu'elle ne prend pas la page en argument. La remise d'attention $w_R$ de
\eqref{eq:distribution} n'est donc pas une propriété de la surface mais de la \textbf{page
servie} : deux fils de composition identique, servis avec des formats différents, n'exposent pas
la même chose.

Cette conclusion est \textbf{retirée} par \S\ref{sec:effetdepage} : les strates ci-dessus tiennent
le rang et les pixels fixes, mais non le \emph{contenu}. Le paragraphe est conservé tel qu'il a été
publié, sa correction venant immédiatement après.

\paragraph{Et une explication de cette note est retirée.} \S\ref{sec:expositionmesuree} attribuait
la forme en deux temps de la courbe mesurée à un \emph{pli d'écran}, et invoquait la colonne
\texttt{serp\_height}. Cette colonne permet de tester l'hypothèse, et elle la réfute : le rang
explique mieux l'examen ($R^2$ de McFadden $0{,}082$) que la hauteur cumulée au-dessus
($0{,}057$), et les pixels n'ajoutent que $0{,}0007$ une fois le rang connu. Le constat --- la
courbe n'est pas une loi de puissance --- tient ; l'explication mécanique tombe.

\subsection{L'effet de page, une fois la composition retirée}
\label{sec:effetdepage}

Le paragraphe précédent comparait les pages enrichies aux autres à rang et à hauteur comparables.
Il lui manquait un contrôle : le \textbf{contenu}. Un encadré de réponse ne s'affiche pas au
hasard --- il répond à une question factuelle, qui n'appelle pas la même lecture qu'une recherche
exploratoire. L'écart mesuré pouvait donc venir du format, ou de ce sur quoi le format apparaît.

\paragraph{Un contrôle qui tranche.} Nous simulons un journal où le format n'a \emph{aucun} effet
sur l'examen, mais où les pages enrichies tombent sur des contenus systématiquement moins
consultés. Stratifier par le rang seul y rend un rapport de risques de $0{,}785$ --- c'est-à-dire
à peu près le chiffre publié plus haut --- alors que la vérité est $1$. Le même protocole apparié
rend $1{,}005$ (IC$_{95}$ $[0{,}993 ; 1{,}016]$), et retrouve $0{,}896$ lorsqu'un effet de
$0{,}90$ est simulé. La stratification par le rang fabrique donc le résultat qu'elle devait
mesurer ; l'appariement, non.

\paragraph{La mesure.} Sur Baidu-ULTR, à \textbf{même document et même rang}, le rapport de
risques sur l'examen vaut $0{,}940$ (IC$_{95}$ $[0{,}916 ; 0{,}964]$, $1\,106$ strates,
$31\,045$ impressions). Un second appariement, par \textbf{requête}, indépendant du premier, rend
$0{,}942$ (IC$_{95}$ $[0{,}869 ; 1{,}022]$) sur sept fois moins d'impressions. Rang par rang,
aucune tendance ne se dégage et tous les intervalles se recouvrent : l'hypothèse la plus simple
qui reste est celle d'un facteur constant.

\paragraph{Ce que cela coûte en pratique.} Un facteur constant est une homothétie, à laquelle un
indice normalisé est insensible ; seul le rang~$1$, qui ne porte jamais rien au-dessus de lui,
empêche l'effet de disparaître tout à fait. Sur quatre mille fils simulés de neuf rangs et quatre
points de vue, ignorer la composition de la page déplace l'indice de $0{,}0025$ en médiane, quand
la convention $1/R$ le déplace de $0{,}0350$ --- \textbf{quatorze fois plus}. Sur la sévérité,
même ordre : $\eta = 0{,}876$ en ignorant la page contre $0{,}854$ sur la page contrefactuelle
nue.

\paragraph{Conséquence.} La remise d'attention $w_R$ de \eqref{eq:distribution} redevient
\textbf{transportable}, ce qui était l'enjeu réel : si elle dépendait de la page, aucun plancher
réglementaire ne pourrait s'écrire sans décrire chaque page servie. La colonne \texttt{format}
du format reste utile --- elle permet de reproduire ce contrôle ailleurs --- mais
elle ne conditionne plus le calcul. Ce qu'il faut exiger d'un journal est d'abord la colonne
d'affichages, sans laquelle l'exposition ne se mesure pas du tout.

\paragraph{Réserves.} L'appariement à contenu fixé ne retient que $1\,106$ strates sur
$404\,478$ : ce sont les documents servis au même rang dans des pages de compositions
différentes, et rien ne garantit qu'ils ressemblent aux autres. L'appariement par requête
concorde mais n'établit rien seul, son intervalle contenant $1$. Enfin, le regroupement de
\texttt{media\_type} en « ordinaire » contre « enrichi » reste un choix de ce travail.

\subsection{L'indice, mesuré sur un fil réel}
\label{sec:indicemesure}

Six sections de cette note répètent la même phrase : \emph{aucun jeu de données public ne porte à
la fois le rang servi et une étiquette de point de vue interprétable}. Elle est fausse, et sa
fausseté est instructive.

\paragraph{Un quatrième journal.} EB-NeRD (\emph{Ekstra Bladet}, RecSys Challenge 2024) donne
pour chaque impression la liste servie \texttt{article\_ids\_inview} et rattache chaque article à
une rubrique déclarée. Les deux colonnes sont là. Le test d'échangeabilité de
\S\ref{sec:exposition} rend pourtant $z = +1{,}05$ ($p = 0{,}29$) sur $232\,887$ fils : l'ordre
enregistré ne porte aucune information de position. Ce n'est pas un défaut de puissance --- la
sévérité que ce journal aurait détectée vaut $0{,}0066$, soit $134$ fois moins que celle mesurée
sur Baidu-ULTR. \texttt{article\_ids\_inview} est un \textbf{ensemble servi}, pas un fil ordonné.

\paragraph{Ce qui devient mesurable.} La forme aveugle au rang, jamais calculée jusqu'ici
ailleurs qu'en simulation. Sur la fenêtre que \S\ref{sec:indice} prescrit --- la
journée-utilisateur --- l'indice vaut $0{,}498$ en médiane, soit $5{,}07$ rubriques également
servies sur $26$, avec $13{,}7\,\%$ des journées sous $0{,}40$ et $50{,}7\,\%$ sous $0{,}50$. Ce
chiffre ne dit pas si c'est peu ou beaucoup ; il donne l'ordre de grandeur qui manquait à toute
discussion de plancher.

\paragraph{Ce qui reste de l'indice exposé.} La composition étant connue et l'ordre non,
l'indice exposé n'est pas déterminé mais \textbf{contraint}. L'énumération exacte de toutes les
mises en ordre distinctes donne un encadrement de largeur médiane $0{,}104$. Trois verdicts
deviennent possibles, et leur asymétrie est le résultat : à un plancher de $0{,}40$, $32{,}4\,\%$
des fils y contreviennent \emph{quelle que soit} leur mise en ordre --- constat opposable sans la
colonne manquante --- $7{,}9\,\%$ sont sûrement conformes, et $59{,}7\,\%$ restent indécidables.
Sans le rang, on peut condamner ; on n'acquitte presque jamais.

\paragraph{Réserves.} \texttt{category\_str} est une \emph{rubrique}, non un point de vue : la
mesure porte sur la diversité thématique exposée, substitut déjà employé sur MIND. La sévérité
$\eta = 0{,}88$ est transportée de Baidu-ULTR, un moteur de recherche ; \S\ref{sec:effetdepage}
a établi qu'elle se transporte d'une page à l'autre, non d'une plateforme à l'autre, d'où un test
de sensibilité qui ne renverse pas la conclusion. L'encadrement exact ne couvre que les fils de
dix contenus ou moins, soit $64\,\%$ du journal.

\subsection{Qui choisit le catalogue choisit l'indice}
\label{sec:catalogue}

\S\ref{sec:indice} porte depuis l'origine une réserve jamais chiffrée : la discrétisation en
points de vue est un choix politique, et qui définit les modalités définit l'indice. Les
journées-utilisateur de \S\ref{sec:indicemesure} permettent enfin de dire \emph{de combien}.

\paragraph{Trois grandeurs, trois comportements.} Le \textbf{niveau} n'a aucun sens absolu : le
même corpus vaut $0{,}498$ sur vingt-six rubriques et $0{,}924$ sur trois. La \textbf{part sous
le plancher} suit : à $0{,}40$, elle passe de $13{,}7\,\%$ à $1{,}5\,\%$ --- le même seuil
devient décoratif sans qu'une seule impression ait changé. L'\textbf{ordre} entre lecteurs, lui,
ne bouge pas : la concordance de rang avec le catalogue complet vaut $1{,}000$ à douze modalités,
$0{,}911$ à six, puis $-0{,}09$ à trois. Il ne se dégrade pas, il s'effondre.

\paragraph{Pourquoi.} Les rubriques rares ne servent presque rien : les onze premières couvrent
$99{,}8\,\%$ du corpus. Regrouper les quinze autres ne fusionne quoi que ce soit que pour une
journée sur mille, d'où un ordre inchangé. Descendre à six force la fusion sur $64{,}5\,\%$ des
journées, à trois sur $94{,}8\,\%$.

\paragraph{La taille ne spécifie pas le catalogue.} Vingt regroupements aléatoires en six classes
rendent une médiane de $0{,}682 \pm 0{,}093$ et une concordance de $0{,}769 \pm 0{,}098$, contre
$0{,}860$ et $0{,}911$ pour le regroupement par fréquence. C'est la \emph{liste} qui doit figurer
dans une norme, non le nombre.

\paragraph{Et la diversité exposée n'existe pas au singulier.} Le corpus porte deux axes
d'étiquetage --- la rubrique thématique et la tonalité déclarée. Leur concordance vaut
$\rho = +0{,}155$ : parmi le quart le plus divers en rubriques, $23{,}8\,\%$ sont aussi dans le
quart le plus divers en tonalité, là où l'indépendance donnerait $25\,\%$. Il faut donc
dire \textbf{sur quel axe} il régule, et ce choix pèse au moins autant que le seuil.

\paragraph{Réserves.} Les catalogues grossiers sont obtenus par regroupement d'un catalogue
existant, ce qui n'est pas un catalogue conçu comme tel. Le corpus est celui d'un seul journal, et
c'est la forme de sa queue de distribution qui explique la stabilité observée jusqu'à douze
modalités.

\subsection{Ce que la contre-expertise a retiré ou restreint}
\label{sec:contreexpertise}

Les instruments de mesure ont été confrontés à la littérature du domaine et à quatre
contre-épreuves. Trois résultats de cette note en sortent modifiés.

\paragraph{Une conclusion retirée.} La comparaison des quatre mesures de diversité
(\S\ref{sec:adverse}) était menée \emph{au même plancher nominal}. Les mesures ne vivant pas sur
la même échelle, cette comparaison n'a pas de sens : à diversité réellement exposée égale, les
deux mesures retenues coûtent la même chose, à moins de $0{,}6\,\%$ de la borne exacte.
L'affirmation selon laquelle la proximité à une cible « résisterait mieux » à l'enterrement est
un artefact d'échelle. Ce qui fait la norme est le \textbf{niveau} exigé et la \textbf{conscience
du rang}, non le choix de la mesure.

\paragraph{Une conclusion restreinte.} Le chiffre « certifiée à $0{,}70$, elle expose $0{,}36$ »
suppose une remise d'attention en $1/R$, c'est-à-dire $\eta = 1$. Recalculé avec la sévérité
\emph{mesurée} d'un bandeau de trois vignettes ($\eta \approx 0{,}1$), le même fil expose
$0{,}747$ : l'enterrement disparaît. À $\eta = 2$ il n'expose que $0{,}157$, et l'échappatoire ne
coûte plus que $2{,}7\,\%$ au lieu de $20{,}8\,\%$. La remise d'attention doit donc être
\textbf{mesurée sur la surface}, et le prix stable est celui de la norme consciente du rang ---
$20{,}8$ à $24{,}1\,\%$ selon la surface.

\paragraph{Une incertitude élargie.} L'estimateur de \S\ref{sec:exposition} suppose un modèle de
clic multiplicatif. La littérature en documente un \emph{affine}, où le lecteur clique par
confiance sur des contenus non pertinents en tête de fil \citep{vardasbi2020}, et démontre que la
pondération par l'inverse de la propension ne peut pas le corriger. Sous ce modèle, la sévérité
estimée dérive de $+12{,}8\,\%$ sans que l'erreur type ne le signale. C'est bien moins grave que
de poser $\eta$ au jugé, mais l'intervalle publié est trop étroit.

\paragraph{Une recommandation confirmée.} L'estimateur doublement robuste fait \emph{moins bien}
que la pondération simple sur l'Open Bandit Dataset ($-4{,}9\,\%$ contre $+2{,}5\,\%$), et la
taille d'échantillon effective ne bouge pas : elle ne dépend que des poids d'importance. Aucun
estimateur ne remplace l'exploration.

\paragraph{Et une mise en garde de la littérature.} Sur Baidu-ULTR --- le jeu même où cette note
mesure $\hat\eta = 1{,}10$ --- \citet{hager2024} constatent que les corrections de biais de
position améliorent la prédiction des clics \emph{sans} améliorer la qualité du classement jugée
par des annotateurs experts. La présence du biais est confirmée par quatre méthodes concordantes,
ce qui corrobore notre chiffre ; l'étape suivante, elle, ne suit pas mécaniquement.

\subsection{Sur le modèle}

\paragraph{Calibration à peine entamée.} Un seul paramètre composite est estimé
$J$, $T$, $\gamma$ et $\alpha$ pris séparément n'ont aucune
procédure d'estimation. La voie la plus prometteuse serait l'inférence des coefficients de
Kramers-Moyal sur des séries d'opinion longues, qui permettrait de \emph{tester} la dérive
plutôt que de l'ajuster.

\paragraph{La seule prédiction propre du modèle n'est pas vérifiée.} Le mécanisme de la charge
émotionnelle $\alpha$ prédit une différence entre registres --- amplification, persistance,
taux de basculement. Quatre mesures indépendantes, dont une réplication sur $440$ sujets et une
annotation en aveugle à double codage, n'en retrouvent aucune trace. L'écart apparent du
corpus pilote était un artefact de sélection.

\paragraph{Une analogie n'est pas une explication.} Rien ici ne démontre que les opinions
\emph{obéissent} à une mécanique statistique ; le travail établit qu'un formalisme emprunté à
la physique en \emph{reproduit} des comportements. Sur l'analogie de départ elle-même, le
verdict est plus net encore.

\paragraph{Les individus ne sont pas des spins.} L'intentionnalité, l'ironie et
l'anticonformisme stratégique n'ont pas d'équivalent dans un modèle de spins ; les opinions
sont multidimensionnelles \citep{deffuant2000}. Surtout, l'environnement n'est pas passif : un
algorithme optimise une fonction de coût, ce qui en fait un acteur stratégique plutôt qu'un
bain thermique. C'est le point le plus fragile de l'analogie --- et c'est aussi ce qui rend
concevable un filtre qui optimise l'indice : une fonction de coût peut être changée.

\paragraph{Portée des simulations.} Les régimes explorés ont été choisis pour leur lisibilité.
Aucune étude systématique de sensibilité n'a été menée, les tailles de systèmes sont modestes,
et les conclusions du modèle retiré étaient qualitatives : elles portaient sur
l'existence de régimes et le sens des dépendances, jamais sur des valeurs numériques
transposables.

\subsection{Sur l'annotation}

\paragraph{Les codeurs ne sont pas indépendants.} Le double codage en aveugle du corpus donne
un $\kappa$ de Fleiss de $0{,}921$, mais les trois codeurs sont des instances du même modèle de
langue : leur accord surestime nécessairement ce que produiraient des juges humains
indépendants. C'est la dernière réserve de méthode qui ne se lève pas par le calcul.

\section{Reproductibilité}

L'intégralité du travail est disponible en accès libre et exécutable en conteneur, avec $556$
tests automatisés. Les vérifications structurantes portent sur la température critique
de la mesure : l'entropie d'un fil dégénéré, la sévérité estimée sur données simulées à
sévérité connue, l'exactitude de l'encadrement de l'indice, et chaque chiffre publié dans
cette note, recalculé depuis les condensés versionnés.

Les journaux d'impressions employés ne sont pas redistribués --- ils pèsent de $135$ Mo à
$2{,}2$ Go et portent chacun sa licence --- mais le dépôt en conserve des condensés vérifiés,
qui rendent à l'identique tous les chiffres publiés ici. Chaque figure est régénérée par un
notebook.

\begin{center}
\url{https://github.com/s-geffroy/Indice-Diversite-Exposee}\\
\url{https://s-geffroy.github.io/Indice-Diversite-Exposee/}
\end{center}

Ce travail est ouvert à la relecture critique. Les retours les plus utiles porteraient sur la
forme retenue de l'indice et sur le niveau d'un éventuel plancher, ou sur les instruments de
mesure des \S\ref{sec:exposition} et \S\ref{sec:evaluation}, où une erreur méthodologique
serait la plus coûteuse.

\bibliography{refs}

\end{document}
